\documentclass[10pt]{beamer}

\usetheme{Madrid}
\usefonttheme{professionalfonts}
\setbeamertemplate{footline}[frame number]
\setbeamercovered{transparent}

\usepackage{algorithm}
\usepackage{algpseudocode}
\usepackage{xcolor}
\renewcommand{\algorithmiccomment}[1]{\hfill$\triangleright$\textcolor{WildStrawberry}{#1}}
\usepackage{float}

\begin{document}

\title{Graphs}
\author{Polina Evgrafova,Catalin Viscun}
\date{November 12, 2025}

\begin{frame}
    \maketitle
\end{frame}

% ------------------------------------------------------------------------------------------------
% GRAPH COLORING SECTION
% ------------------------------------------------------------------------------------------------

\section{Graph Coloring Basics}

\begin{frame}{Grundidee der Graphfärbung}
    \small
    \begin{itemize}
        \item Das Ziel der Graphfärbung ist es, jedem Knoten eines Graphen eine Farbe zuzuweisen, sodass keine zwei benachbarten Knoten dieselbe Farbe erhalten.
        \item Die kleinste Anzahl an Farben, die dafür nötig ist, wird \textbf{chromatische Zahl} $\chi(G)$ genannt.
        \item Dieses Prinzip dient als Modell für zahlreiche Planungs- und Zuweisungsprobleme, bei denen sich bestimmte Ereignisse, Ressourcen oder Entitäten nicht überlappen dürfen.
        \item In der Stundenplanerstellung entspricht eine Farbe einem \textbf{Zeitslot}: zwei verbundene Knoten dürfen nicht gleichzeitig stattfinden.
        \item Die Färbung macht Konflikte mathematisch sichtbar und ermöglicht algorithmische Optimierung.
    \end{itemize}

    \vspace{1em}
    \textbf{Formale Definition:}
    \[
        c: V(G) \rightarrow \{1, 2, \dots, k\} \quad \text{mit} \quad (u,v) \in E(G) \Rightarrow c(u) \neq c(v)
    \]
\end{frame}

\begin{frame}{Algorithmische Grundidee}
    \small
    \begin{itemize}
        \item Das Problem ist \textbf{NP-schwer}: Es gibt keine bekannte polynomielle Methode, die garantiert die minimale Anzahl an Farben findet.
        \item Praktisch werden daher heuristische oder approximative Verfahren genutzt, die in endlicher Zeit brauchbare Ergebnisse liefern.
        \item Einfache Heuristik: \textbf{Greedy Coloring}
        \begin{itemize}
            \item Knoten werden nach bestimmtem Kriterium (z. B. Knotengrad) sortiert.
            \item Jeder Knoten erhält die kleinste mögliche Farbe, die nicht bei seinen Nachbarn vorkommt.
        \end{itemize}
        \item Erweiterte Verfahren: \textbf{DSATUR} (Degree of Saturation) priorisiert Knoten, deren Nachbarschaft bereits stark eingefärbt ist.
        \item Für exakte Lösungen werden \textbf{Backtracking-} oder \textbf{Branch\&Bound-Methoden} verwendet, die jedoch exponentiell wachsen.
    \end{itemize}
\end{frame}

\begin{frame}{Mathematische Modellierung (ILP)}
    \small
    \[
        x_{ij} =
        \begin{cases}
            1 & \text{wenn Knoten } i \text{ Farbe } j \text{ erhält} \\
            0 & \text{sonst}
        \end{cases}
    \]
    Nebenbedingungen:
    \[
        \sum_j x_{ij} = 1 \quad \forall i
    \]
    \[
        x_{ij} + x_{kj} \le 1 \quad \forall (i,k) \in E(G)
    \]
    Ziel:
    \[
        \min \sum_j y_j \quad \text{mit } y_j \ge x_{ij}
    \]

    \vspace{0.5em}
    \begin{itemize}
        \item Diese ILP-Formulierung erlaubt es, Graphfärbung als lineares Optimierungsproblem zu behandeln.
        \item Das Modell wird häufig in Stundenplanungssoftware verwendet, da es leicht mit weiteren Nebenbedingungen kombinierbar ist (Räume, Präferenzen, Lehrer, Ressourcen).
    \end{itemize}
\end{frame}

% ------------------------------------------------------------------------------------------------
% SPECIAL CASES
% ------------------------------------------------------------------------------------------------

\section{Spezialfälle der Graphfärbung}

\begin{frame}{Vertex Coloring}
    \small
    \begin{itemize}
        \item Klassischste Form der Graphfärbung.
        \item Knoten repräsentieren Entitäten (z. B. Kurse oder Prüfungen), Kanten stellen Konflikte dar.
        \item Farben stehen für Zeitslots, sodass keine zwei verbundene Knoten gleichzeitig stattfinden.
        \item Wird in der Stundenplanung eingesetzt, um die minimale Anzahl an Zeitslots zu bestimmen.
        \item Effizient kombinierbar mit zusätzlichen Restriktionen wie Raumkapazität oder Dozentenverfügbarkeit.
    \end{itemize}
\end{frame}

\begin{frame}{Edge Coloring}
    \small
    \begin{itemize}
        \item Bei der Kantenfärbung werden nicht die Knoten, sondern die \textbf{Kanten} eingefärbt.
        \item Zwei Kanten, die denselben Knoten teilen, dürfen nicht dieselbe Farbe haben.
        \item Farben können hier für Zeitslots stehen, während Knoten Ressourcen (z. B. Räume, Lehrer) darstellen.
        \item Ziel: Jede Ressource wird in einem bestimmten Zeitfenster nur einmal genutzt.
        \item Besonders geeignet für Probleme, bei denen Ressourcenverteilung zentral ist.
    \end{itemize}
\end{frame}

\begin{frame}{List Coloring}
    \small
    \begin{itemize}
        \item Jeder Knoten besitzt eine eingeschränkte Menge möglicher Farben (eine individuelle „Liste“).
        \item Modelliert Verfügbarkeit oder Präferenzen, z. B.:
        \begin{itemize}
            \item Eine Lehrperson kann nur Montag und Mittwoch unterrichten.
            \item Ein Raum steht nur zu bestimmten Zeiten zur Verfügung.
        \end{itemize}
        \item Aufgabe: Finde eine gültige Färbung, bei der die Farbe jedes Knotens aus seiner Liste stammt.
        \item Erweitert klassische Färbung um reale Einschränkungen.
    \end{itemize}
\end{frame}

\begin{frame}{Multi-Coloring}
    \small
    \begin{itemize}
        \item Ein Knoten erhält nicht nur eine, sondern mehrere Farben.
        \item Repräsentiert Entitäten, die mehrfach pro Woche oder mehrfach im Plan auftreten müssen (z. B. drei Vorlesungstermine).
        \item Bedingung bleibt: Keine Nachbarknoten dürfen eine gemeinsame Farbe besitzen.
        \item Praktische Bedeutung in Stundenplänen, die Mehrfachsitzungen pro Kurs beinhalten.
    \end{itemize}
\end{frame}

\begin{frame}{Weighted und Dynamic Coloring}
    \small
    \textbf{Weighted Coloring}
    \begin{itemize}
        \item Konflikte oder Präferenzen werden mit Gewichten versehen.
        \item Ziel: Minimierung der Gesamtkosten oder Strafen bei Verletzungen weicher Regeln.
        \item Beispiel: Ein unerwünschter, aber zulässiger Paralleltermin hat geringes Gewicht, ein verbotener Konflikt hohes Gewicht.
    \end{itemize}

    \vspace{0.5em}
    \textbf{Dynamic Coloring}
    \begin{itemize}
        \item Verwendet für Situationen, in denen sich der Graph im Lauf der Zeit verändert (z. B. neue Kurse, Raumänderungen).
        \item Ziel: Anpassung der bestehenden Färbung mit minimalen Änderungen.
        \item Relevant bei inkrementeller Stundenplanaktualisierung.
    \end{itemize}
\end{frame}

\begin{frame}{Hypergraph Coloring}
    \small
    \begin{itemize}
        \item Erweiterung auf \textbf{Hypergraphen}, bei denen Kanten mehr als zwei Knoten verbinden können.
        \item Eine Hyperkante kann z. B. alle Veranstaltungen umfassen, die denselben Raum gleichzeitig beanspruchen.
        \item Färbung stellt sicher, dass kein Element einer Hyperkante dieselbe Farbe hat wie ein anderes derselben Gruppe.
        \item Wird eingesetzt, wenn Abhängigkeiten nicht nur paarweise auftreten (z. B. Team-Teaching oder Gruppenkurse).
    \end{itemize}
\end{frame}

\begin{frame}{Fazit}
    \small
    \begin{itemize}
        \item Graph Coloring ist das grundlegende Prinzip graphbasierter Stundenplanung.
        \item Alle Varianten (Vertex-, Edge-, List-, Multi-, Weighted-, Hypergraph Coloring) erweitern das Grundmodell um spezifische Randbedingungen.
        \item Effektiv, weil Konflikte explizit abgebildet werden und algorithmisch lösbar sind.
        \item Für mittlere Problemgrößen liefern heuristische Verfahren meist in Sekunden bis Minuten brauchbare, konfliktfreie Lösungen.
        \item In großen, realen Szenarien oft Teil eines hybriden Systems mit Constraint Programming oder Linearer Optimierung.
    \end{itemize}
\end{frame}


\end{document}
